\documentclass[11pt]{article}

\usepackage[margin=1in]{geometry}
\usepackage{amsmath,amssymb}
\usepackage{bm}

\title{Finite-size scaling with GP scaling functions\\
and FSS mixture discrepancy ($\kappa$, $\omega$)}
\author{}
\date{}

\begin{document}
\maketitle

\section{Overview}

Extracting critical exponents from simulations of statistical field
theories is traditionally done with a complicated, but ad~hoc,
frequentist procedure.
The goal of this project, in the spirit of~\cite{harada}, is to
systematize that procedure using Bayesian probability.

\section{Finite-size scaling from the renormalization group}

View renormalization as an operator $R_b : K \to K$ (for $b\in\mathbb{R}$)
on the coupling space of a theory $K$.
We want the eigenvalues of the linearization $J_b$ of $R_b$ around a
fixed point of interest.
If
\begin{equation}
  J_b(u_i) = b^{y_i}\, u_i,
\end{equation}
then $y_i$ is the RG eigenvalue of the corresponding scaling field $u_i$
(often written $b^{-x_i}$ for the dual operator scaling dimensions; we
keep the $y_i$ convention below).

For an observable $O$, define the pushforward
$R_b^* O = O\circ R_b$.
This likewise has an eigendecomposition.
For an eigenoperator $\Phi$ of scaling dimension $x_O$,
\begin{equation}
  \langle\Phi\rangle(t,\{u_i\})
    = b^{-x_O}\,
      \langle\Phi\rangle\!\bigl(b^{y_t} t,\, \{b^{y_i} u_i\}\bigr).
\end{equation}
Taking the finite-size identification $b = L$,
\begin{equation}
  \label{eq:fss-eigen}
  \langle\Phi\rangle(t,\{u_i\})
    = L^{-x_O}\,
      \langle\Phi\rangle\!\bigl(t L^{y_t},\, \{u_i L^{y_i}\}\bigr).
\end{equation}
A general observable is a linear combination of such eigenoperators, so
\begin{equation}
  \label{eq:fss-general}
  \boxed{
  \langle O\rangle_L
    = L^{-x_0}
      \Biggl[
        F_0\bigl(\{u_i L^{y_i}\}\bigr)
        +
        \sum_{a>0}
        L^{-(x_a-x_0)}
        F_a\bigl(\{u_i L^{y_i}\}\bigr)
      \Biggr]
  }
\end{equation}
where $x_0$ is the leading (most relevant) operator dimension contributing
to $O$, and the sum collects subleading operators.

\subsection{Ising magnetization moments}

Specialize to the 2D Ising universality class with a single relevant
thermal perturbation.
Write
\begin{equation}
  t = \frac{T-T_c}{T_c}
\end{equation}
for the reduced temperature, so the thermal scaling field is
$u_t \propto t$ with RG eigenvalue
\begin{equation}
  y_t = \frac{1}{\nu}.
\end{equation}
The leading irrelevant bulk operator has eigenvalue $y_{\mathrm{irr}} = -\omega$
(Wegner's correction-to-scaling exponent), and we write $u$ for its
(nonlinear) scaling field.
Keeping only these two fields, the arguments of the scaling functions are
\begin{equation}
  z := t\, L^{1/\nu},
  \qquad
  u\, L^{-\omega}.
\end{equation}

For the absolute magnetization density $m = \langle |M|\rangle / N$,
the leading magnetic scaling dimension is $x_0 = \beta/\nu$, so
\begin{equation}
  \label{eq:m-fss}
  m(t,L)
    = L^{-\beta/\nu}\,
      \Bigl[
        F_m\!\bigl(z,\, u L^{-\omega}\bigr)
        + O\!\bigl(L^{-(x_1-\beta/\nu)}\bigr)
      \Bigr].
\end{equation}
The moments $\langle m^2\rangle$ and $\langle m^4\rangle$ are controlled by
$2$ and $4$ powers of the magnetic operator, hence
\begin{align}
  \label{eq:m2-fss}
  m^2(t,L)
    &= L^{-2\beta/\nu}\,
      \Bigl[
        F_{m^2}\!\bigl(z,\, u L^{-\omega}\bigr)
        + \cdots
      \Bigr],
  \\
  \label{eq:m4-fss}
  m^4(t,L)
    &= L^{-4\beta/\nu}\,
      \Bigl[
        F_{m^4}\!\bigl(z,\, u L^{-\omega}\bigr)
        + \cdots
      \Bigr].
\end{align}
(Equivalently: collapsing with $L^{k\beta/\nu}$ for the $k$-th moment removes
the leading $L$-power and leaves a universal function of
$(z,\, u L^{-\omega})$.)

Expanding the leading scaling function in the irrelevant direction near
the fixed point,
\begin{equation}
  \label{eq:wegner-expand}
  F_c(z,\, u L^{-\omega})
    = f_c(z) + (u L^{-\omega})\, g_c(z) + O(L^{-2\omega}),
\end{equation}
so the collapsed fields
\begin{equation}
  \Phi_m := m\, L^{\beta/\nu},
  \qquad
  \Phi_{m^2} := m^2\, L^{2\beta/\nu},
  \qquad
  \Phi_{m^4} := m^4\, L^{4\beta/\nu}
\end{equation}
satisfy, to leading correction,
\begin{equation}
  \label{eq:phi-wegner}
  \Phi_c(z,L)
    = f_c(z) + u\, L^{-\omega}\, g_c(z) + \cdots.
\end{equation}
Thus the natural finite-size correction to a collapsed observable is
multiplicative in $L^{-\omega}$: that is why the Wegner factor appears as
a power of $L$, not of $t$ alone.

\subsection{Why $a = L^{-\omega} + \kappa\,|t|^{\omega\nu}$}

Equation~\eqref{eq:phi-wegner} already motivates an amplitude proportional
to $L^{-\omega}$ at $t=0$.
Away from criticality the irrelevant scaling field itself depends on $t$,
and bulk Wegner corrections produce terms $\sim |t|^{\Delta}$ with
$\Delta = \omega\nu$.
A minimal two-term ansatz that (i)~matches the finite-size correction
$L^{-\omega}$ and (ii)~matches the bulk correction $|t|^{\omega\nu}$, while
remaining homogeneous under the collapse $(t,L)\mapsto z$, is
\begin{equation}
  \label{eq:a-fss}
  a(t,L)
    = L^{-\omega} + \kappa\, |t|^{\omega\nu}
    = L^{-\omega}\bigl(1 + \kappa\, |z|^{\omega\nu}\bigr).
\end{equation}
The factored form makes the RG content explicit: the whole correction is
$O(L^{-\omega})$ times a smooth function of the collapse coordinate $z$,
and $1+\kappa|z|^{\omega\nu}$ is a simple positive surrogate for that
$z$-dependence (with relative weight $\kappa$).
In particular:
\begin{itemize}
  \item $L^{-\omega}$ comes from the RG eigenvalue $y_{\mathrm{irr}}=-\omega$ of
    the leading irrelevant operator in~\eqref{eq:fss-general}
    and~\eqref{eq:wegner-expand};
  \item $|t|^{\omega\nu}$ is the bulk conjugate
    ($|t|^{\Delta}$ with $\Delta=\omega\nu$), equivalently
    $|z|^{\omega\nu} L^{-\omega}$ after collapse;
  \item $a\to 0$ whenever $t\to 0$ and $L\to\infty$ together, so the
    asymptotic scaling window is recovered.
\end{itemize}

\section{The Bayesian setup}

As in any Bayesian problem, we start from unknowns and knowns.
Concretely, for the 2D Ising model with one relevant thermal
perturbation:

\paragraph{Unknowns.}
The critical parameters $\theta = (T_c,\nu,\beta)$ (and the correction
exponents / amplitudes $\omega,\kappa$), together with the universal
scaling functions $f_c$ (and discrepancy functions $g_c$) for each
observable channel $c\in\{m,m^2,m^4\}$.

\paragraph{Knowns.}
\begin{itemize}
  \item Critical exponents are positive; $T_c$ is positive.
  \item Scaling functions are smooth, with a characteristic length scale
    in $z$ that we encode through a GP kernel.
  \item We can estimate observables at any $(T,L)$ by Monte Carlo, up to
    a statistical uncertainty $\sigma$; reducing $\sigma$ is costly and
    compute is finite.
  \item For fixed $\theta$, the collapsed observables $\Phi_c$ become
    increasingly accurate descriptions of a universal $f_c(z)$ in the
    limit $t\to 0$, $L\to\infty$ with $z = t L^{1/\nu}$ fixed
    (Eqs.~\eqref{eq:fss-general}--\eqref{eq:phi-wegner}).
\end{itemize}

The data are Monte Carlo estimates
\[
  D
  = \bigl\{
      (L_i,T_i,\, m_i,m^2_i,m^4_i,\,
       \sigma_{m,i},\sigma_{m^2,i},\sigma_{m^4,i})
    \bigr\}_{i=1}^{n}.
\]
For each observation form the collapsed fields and propagate the MC
errors,
\begin{align}
  \Phi_{m,i}
    &= m_i\, L_i^{\beta/\nu},
  &
  \sigma_{\Phi_m,i}
    &= \sigma_{m,i}\, L_i^{\beta/\nu},
  \\
  \Phi_{m^2,i}
    &= m^2_i\, L_i^{2\beta/\nu},
  &
  \sigma_{\Phi_{m^2},i}
    &= \sigma_{m^2,i}\, L_i^{2\beta/\nu},
  \\
  \Phi_{m^4,i}
    &= m^4_i\, L_i^{4\beta/\nu},
  &
  \sigma_{\Phi_{m^4},i}
    &= \sigma_{m^4,i}\, L_i^{4\beta/\nu}.
\end{align}
Write $\Phi^{(c)} = (\Phi_{c,1},\ldots,\Phi_{c,n})^\top$ for channel
$c\in\{m,m^2,m^4\}$, and
\begin{equation}
  t_i = \frac{T_i - T_c}{T_c},
  \qquad
  z_i = t_i\, L_i^{1/\nu}.
\end{equation}

\section{Latent model: universal GP + FSS mixture discrepancy}

Rather than fitting parametric forms for $f_c$ and $g_c$, place
independent zero-mean GP priors on $z$,
\begin{equation}
  f^{(c)} \sim \mathrm{GP}(0,\, k_{\ell_f,\eta_f}),
  \qquad
  g^{(c)} \sim \mathrm{GP}(0,\, k_{\ell_g,\sigma_g}),
\end{equation}
with RBF kernel
\begin{equation}
  k_{\ell,\eta}(z,z')
    = \eta^2
      \exp\!\Bigl(-\frac{(z-z')^2}{2\ell^2}\Bigr).
\end{equation}
Using the amplitude~\eqref{eq:a-fss}, define the unit-interval mixture gate
\begin{equation}
  \label{eq:pi-gate}
  \pi(t,L)
    = \frac{a(t,L)}{1 + a(t,L)}
    \in [0,1).
\end{equation}
The collapsed observation model is the \emph{mixture} coupling
\begin{equation}
  \label{eq:phi-model}
  \Phi^{(c)}_i
    = \bigl(1-\pi_i\bigr)\, f^{(c)}(z_i)
      + \pi_i\, g^{(c)}(z_i)
      + \varepsilon^{(c)}_i,
  \qquad
  \varepsilon^{(c)}_i
    \sim \mathcal{N}\!\bigl(0,\, \sigma_{\Phi_c,i}^2\bigr),
\end{equation}
where $\pi_i = \pi(t_i,L_i)$.
Channels share $(\omega,\kappa,\ell_f,\eta_f,\ell_g,\sigma_g)$ but have
independent draws of $(f^{(c)},g^{(c)})$.
There is no extra $L^{k\beta/\nu}$ factor on $a$ in $\Phi$-space: after
collapsing,~\eqref{eq:a-fss} already multiplies the correction channel.

As $t\to 0$ and $L\to\infty$, $a\to 0$ so $\pi\to 0$
and~\eqref{eq:phi-model} recovers the plain universal model
$\Phi = f + \varepsilon$.
When $a$ is large (small $L$ and/or large $|z|$), $\pi\to 1$ and the point
is explained by $g$ instead of $f$: out-of-window observations do not train
the universal scaling function.
When $\kappa=0$,~\eqref{eq:a-fss} reduces to pure $L^{-\omega}$ gating,
\begin{equation}
  \pi = \frac{L^{-\omega}}{1+L^{-\omega}},
\end{equation}
the mixture analogue of the leading Wegner correction
in~\eqref{eq:phi-wegner}.

\section{Prior}

Collect the inferred scalars into
\begin{equation}
  \psi
    = \bigl(
        T_c,\, \nu,\, \beta,\,
        \omega,\, \kappa,\,
        \ell_f,\, \eta_f,\, \ell_g,\, \sigma_g
      \bigr)
\end{equation}
(components that are held fixed are simply omitted from $\psi$).
The prior factorises:
\begin{align}
  p(\psi)
    &= p(T_c)\, p(\nu)\, p(\beta)
  \nonumber\\
    &\qquad
      \times p(\omega)\, p(\kappa)
  \nonumber\\
    &\qquad
      \times p(\ell_f)\, p(\eta_f)\, p(\ell_g)\, p(\sigma_g).
\end{align}
Explicitly (matching the defaults used in the code),
\begin{align}
  T_c
    &\sim \mathrm{Uniform}(2.0,\, 2.5),
  \\
  \nu
    &\sim \mathrm{Uniform}(0.5,\, 1.5),
  \\
  \beta
    &\sim \mathrm{Uniform}(0.05,\, 0.25),
  \\
  \omega
    &\sim \mathrm{Uniform}(1,\, 10),
  \\
  \kappa
    &\sim \mathrm{Uniform}(0,\, 20),
  \\
  \ell_f
    &\sim \mathrm{LogNormal}\bigl(\log \ell_f^{\star},\, \sigma_\ell\bigr),
  \\
  \eta_f
    &\sim \mathrm{HalfNormal}(\sigma_{\eta_f}),
  \\
  \ell_g
    &\sim \mathrm{LogNormal}\bigl(\log \ell_g^{\star},\, \sigma_\ell\bigr)
      \quad\text{(or tied to $\ell_f$)},
  \\
  \sigma_g
    &\sim \mathrm{HalfNormal}(1),
\end{align}
where $\ell_f^{\star},\ell_g^{\star}$ are data-dependent length-scale
guesses (fractions of the provisional $z$-span) and
$\sigma_\ell,\sigma_{\eta_f}$ are fixed prior scales.
Independently of $\psi$, the latent fields have GP priors
\begin{equation}
  p\bigl(\{f^{(c)},g^{(c)}\}_c \bigm| \psi\bigr)
    = \prod_{c\in\{m,m^2,m^4\}}
      \mathrm{GP}(f^{(c)};\ell_f,\eta_f)\,
      \mathrm{GP}(g^{(c)};\ell_g,\sigma_g).
\end{equation}

\section{Likelihood}

Conditional on the latents and $\psi$, channels are independent and
Gaussian:
\begin{equation}
  p\bigl(D \bigm| \{f^{(c)},g^{(c)}\}_c,\, \psi\bigr)
    = \prod_{c\in\{m,m^2,m^4\}}
      \prod_{i=1}^{n}
      \mathcal{N}\!\bigl(
        \Phi^{(c)}_i \;\big|\;
        (1-\pi_i)\, f^{(c)}(z_i) + \pi_i\, g^{(c)}(z_i),\;
        \sigma_{\Phi_c,i}^2
      \bigr).
\end{equation}
Integrating out $(f^{(c)},g^{(c)})$ analytically gives the marginal
likelihood used for inference:
\begin{equation}
  p(D\mid\psi)
    = \prod_{c\in\{m,m^2,m^4\}}
      \mathcal{N}\!\bigl(\Phi^{(c)} \;\big|\; 0,\, K^{(c)}(\psi)\bigr),
\end{equation}
with channel covariances
\begin{equation}
  \label{eq:K-mixture}
  K^{(c)}_{ij}(\psi)
    = (1-\pi_i)(1-\pi_j)\, k_{\ell_f,\eta_f}(z_i,z_j)
      + \pi_i\pi_j\, k_{\ell_g,\sigma_g}(z_i,z_j)
      + \sigma_{\Phi_c,i}^2\, \delta_{ij},
\end{equation}
where $\pi_i = \pi(t_i,L_i)$ from~\eqref{eq:pi-gate}.
Writing this out,
\begin{align}
  \log p(D\mid\psi)
    &= \sum_{c\in\{m,m^2,m^4\}}
      \Biggl[
        -\tfrac12\, {\Phi^{(c)}}^\top
          \bigl(K^{(c)}\bigr)^{-1} \Phi^{(c)}
  \nonumber\\
    &\qquad\qquad
        -\tfrac12\log\det K^{(c)}
        -\tfrac{n}{2}\log(2\pi)
      \Biggr].
\end{align}
(If a channel is disabled, omit its term.)

\section{Posterior}

Bayes' rule then gives
\begin{equation}
  p(\psi\mid D)
    \propto
    p(\psi)\, p(D\mid\psi).
\end{equation}
Sampling targets $\psi$ with the latent GPs integrated out; predictions
for $f^{(c)}$ and $g^{(c)}$ are recovered from the usual GP posterior
conditional on each draw of $\psi$.

\begin{thebibliography}{9}

\bibitem{harada}
  K.~Harada,
  \emph{Bayesian inference of the critical exponents from finite-size
  scaling analysis},
  Phys.\ Rev.\ E \textbf{84}, 056704 (2011).

\end{thebibliography}

\end{document}
